% Einheit 3 von 4 – Normalform und p-q-Formel (bank/quadratische-gleichungen/e3.jsonl, zusammenbau v0.1)
%% TODO Zweigzeile Teil 1: was hier gelernt wird (Satz in Schülersprache aus dem Einheitstitel) – Einheit 3
\einheitenkopf[e3]{Einheit 3 von 4 · Normalform und p-q-Formel}
\zweigzeile{ab Kl. 9, je nach Buch bis Kl. 10 · P10 oft}
\setcounter{aufgabe}{15}

%% TODO Titel: Ich-kann-Satz fehlt in der Bank (Platzhalter: Kettenname) – Nr. 16
%% TODO Anweisung: ein Satz über der ersten Teilaufgabe fehlt in der Bank – Nr. 16
\begin{aufgabe}{Welche Form?}
\begin{teile}
\teil $x^2 - 3x + 1 = 0$ – welche Form, welcher Weg?\\ \kreuz{reinquadratisch – Wurzelziehen} \\ \kreuz{Produkt gleich null – Nullprodukt} \\ \kreuz{quadrierte Klammer – Rückwärtsrechnen} \\ \kreuz{Normalform – Formel}
\end{teile}
\end{aufgabe}

%% TODO Titel: Ich-kann-Satz fehlt in der Bank (Platzhalter: Kettenname) – Nr. 17
%% TODO Anweisung: ein Satz über der ersten Teilaufgabe fehlt in der Bank – Nr. 17
\begin{aufgabe}{Ist die Vorzahl von x² eins?}
\begin{teile}
\teil $4x^2 + 8x - 12 = 0$ – vor der Formel teilen?\\ \kreuz{nicht teilen} \\ \kreuz{durch $4$} \\ \kreuz{durch $8$}
\end{teile}
\end{aufgabe}

%% TODO Titel: Ich-kann-Satz fehlt in der Bank (Platzhalter: Kettenname) – Nr. 18
%% TODO Anweisung: ein Satz über der ersten Teilaufgabe fehlt in der Bank – Nr. 18
\begin{aufgabe}{p-q-Formel}
%% TODO \rechenplatz{n}: Schrittzahl des Grundfalls fehlt in der Bank (nur bei fester Schreibform, 2.3 b)
\begin{teile}
\teil $x^2 - x - 12 = 0$ – p und q? p = \leerfeld , q = \leerfeld
\end{teile}
\begin{gleichungsraster}[2]
\gl{\text{x^2 + 6x + 8 = 0}} &
\gl{\text{x^2 - 4x - 12 = 0}} \\
\gl{\text{x^2 + 5 = 6x}} &
\gl{\text{5x^2 - 20x + 15 = 0}} \\
\gl{\text{x^2 - 3x + 2 = 0}} &
\gl{\text{x^2 - 10x + 25 = 0}} \\
\gl{\text{x^2 + 2x + 5 = 0}} &
\end{gleichungsraster}
\begin{teile}
\teil $x^2 - 8x + 7 = 0$ – Diskriminante $D = \left(\frac{p}{2}\right)^2 - q$ und Zahl der Lösungen? D = \leerfeld ; Lösungen: \leerfeld
\end{teile}
\begin{gleichungsraster}[2]
\gl{\text{x^2 - 6x + 4 = 0}} &
\end{gleichungsraster}
\begin{teile}
\teil Probe: Ist $x = -2$ Lösung von $x^2 - 3x - 10 = 0$? \janein
\end{teile}
\begin{gleichungsraster}[2]
\gl{\text{(x + 2)^2 = 5x + 10}} &
\end{gleichungsraster}
\begin{teile}
\teil $x^2 - 12x = 0$ – welcher Weg, welche Lösungen? Weg: \leerfeld ; x1 = \leerfeld , x2 = \leerfeld
\teil Die Funktion $f$ mit $f(x) = 2x^2 - 2x - 40$ – Nullstellen? \hfill (P10 2020 OS) \\ x1 = \leerfeld , x2 = \leerfeld
\end{teile}
\end{aufgabe}

%% TODO Titel: Ich-kann-Satz fehlt in der Bank (Platzhalter: Kettenname) – Nr. 19
%% TODO Anweisung: ein Satz über der ersten Teilaufgabe fehlt in der Bank – Nr. 19
\begin{aufgabe}{Gerade und Parabel gleichsetzen als Anwendung}
\begin{teile}
\teil Parabel $y = x^2$ und Gerade $y = x + 6$ – an welchen Stellen $x$ schneiden sie sich? x1 = \leerfeld , x2 = \leerfeld
\end{teile}
\end{aufgabe}

%% TODO Titel: Ich-kann-Satz fehlt in der Bank (Platzhalter: Kettenname) – Nr. 20
%% TODO Anweisung: ein Satz über der ersten Teilaufgabe fehlt in der Bank – Nr. 20
\begin{aufgabe}{p-q-Formel – Fehler finden, Begründen, Darstellungswechsel, Anwendung}
\begin{teile}
\teil Wo ist der Fehler? \rechnung{3x^2 + 6x - 24 &= 0 \\ x_{1,2} &= -3 \pm \sqrt{9 + 24}}
\teil Begründe, warum du bei $3x^2 + 9x - 12 = 0$ vor der Formel durch $3$ teilen musst.
\teil Die Parabel $y = x^2 - 4x + 3$ ist gezeichnet. Lösungen von $x^2 - 4x + 3 = 0$ am Graphen?

\begin{ksys}[xmin=-2,xmax=6,ymin=-2,ymax=5,ablesen] \parabel{1}{2}{-1}{} \end{ksys}
x1 = \leerfeld , x2 = \leerfeld
\teil Ein Ball wird geworfen. Für seine Höhe gilt $h = -0{,}5x^2 + 3x + 8$ ($x$ waagerechte Entfernung in m, $h$ Höhe in m). In welcher Entfernung trifft er auf den Boden? \leerfeld[m]
\end{teile}
\end{aufgabe}
